\section{AHSME 1977}
        \begin{enumerate}
        	\item (\textit{AHSME 1977}) If $y = 2x$ and $z = 2y$, then $x + y + z$ equals 


 $\text{(A)}\ x \qquad  \text{(B)}\ 3x \qquad  \text{(C)}\ 5x \qquad  \text{(D)}\ 7x \qquad  \text{(E)}\ 9x$


 
 \textbf{Answer: }D



\textbf{Solution}

Solution by e\_power\_pi\_times\_i

 $x+y+z = x+(2x)+(4x) = \boxed{\text{(D)}\ 7x}$

 


	\item (\textit{AHSME 1977}) Which one of the following statements is false? All equilateral triangles are


 $\textbf{(A)}\ \text{ equiangular}\qquad \textbf{(B)}\ \text{isosceles}\qquad \textbf{(C)}\ \text{regular polygons }\qquad\\ \textbf{(D)}\ \text{congruent to each other}\qquad \textbf{(E)}\ \text{similar to each other}$


 
 \textbf{Answer: }D



\textbf{Solution}

Solution by e\_power\_pi\_times\_i

 Equilateral triangles can be in different sizes, therefore they are not $\boxed{\textbf{(D)}\ \text{congruent to each other}}$

 


	\item (\textit{AHSME 1977}) A man has $2.73 in pennies, nickels, dimes, quarters and half dollars. If he has an equal number of coins of each kind, then the total number of coins he has is 


 $\text{(A)}\ 3 \qquad  \text{(B)}\ 5 \qquad  \text{(C)}\ 9 \qquad  \text{(D)}\ 10 \qquad  \text{(E)}\ 15$


 
 \textbf{Answer: }E



\textbf{Solution}

Solution by e\_power\_pi\_times\_i

 Denote the number of pennies, nickels, dimes, quarters, and half-dollars by $x$. Then $x+5x+10x+25x+50x = 273$, and $x = 3$. The total amount of coins is $5x$, which is $\textbf{\text{(E)}}\ 15$.

 


	\item (\textit{AHSME 1977}) 
                    \begin{figure}[H]
                        centering
                        \includegraphics[width=0.5\textwidth]{Problem/AHSME_I_1977_4_Problem_Figure_1.png}
                    \end{figure}
                    


 In triangle $ABC, AB=AC$ and $\measuredangle A=80^\circ$. If points $D, E$, and $F$ lie on sides $BC, AC$ and $AB$, respectively, and $CE=CD$ and $BF=BD$, then $\measuredangle EDF$ equals


 $\textbf{(A) }30^\circ\qquad \textbf{(B) }40^\circ\qquad \textbf{(C) }50^\circ\qquad \textbf{(D) }65^\circ\qquad \textbf{(E) }\text{none of these}$


 
 \textbf{Answer: }C



\textbf{Solution}

Solution by e\_power\_pi\_times\_i

 Because $\measuredangle A=80^\circ$, $\measuredangle B=\measuredangle C=50^\circ$. Then, because triangles $CDE$ and $BDF$ are isosceles, $\measuredangle CDE=\measuredangle BDF=65^\circ$. $\measuredangle EDF=180^\circ - 2(65^\circ)=\textbf{(C) }50^\circ$

 


	\item (\textit{AHSME 1977}) The set of all points $P$ such that the sum of the (undirected) distances from $P$ to two fixed points $A$ and $B$ equals the distance between $A$ and $B$ is


 $\textbf{(A) }\text{the line segment from }A\text{ to }B\qquad \textbf{(B) }\text{the line passing through }A\text{ and }B\qquad\\ \textbf{(C) }\text{the perpendicular bisector of the line segment from }A\text{ to }B\qquad\\ \textbf{(D) }\text{an ellipse having positive area}\qquad \textbf{(E) }\text{a parabola}$


 
 \textbf{Answer: }A



\textbf{Solution}

Solution by e\_power\_pi\_times\_i

 The answer is $\textbf{(A)}$ because $P$ has to be on the line segment $AB$ in order to satisfy $PA+PB=AB$.

 


	\item (\textit{AHSME 1977}) If $x, y$ and $2x + \frac{y}{2}$ are not zero, then
$\left( 2x + \frac{y}{2} \right)^{-1} \left[(2x)^{-1} + \left( \frac{y}{2} \right)^{-1} \right]$
equals


 $\textbf{(A) }1\qquad \textbf{(B) }xy^{-1}\qquad \textbf{(C) }x^{-1}y\qquad \textbf{(D) }(xy)^{-1}\qquad  \textbf{(E) }\text{none of these}$


 
 \textbf{Answer: }D



\textbf{Solution}

We can write $\left( 2x+ \frac{y}{2} \right)^{-1} = \left( \frac{4x+y}{2} \right)^{-1} = \frac{2}{4x+y}$.

 Then, the expression simplifies to 
 \[\left( 2x + \frac{y}{2} \right)^{-1} \left[(2x)^{-1} + \left( \frac{y}{2} \right)^{-1} \right] \Rightarrow \frac{2}{4x+y} \left( \frac{1}{2x} + \frac{2}{y} \right) \Rightarrow \frac{2}{4x+y} \left( \frac{y+4x}{2xy} \right) \Rightarrow \frac{1}{xy}.\]

 
Thus, our answer is $\textbf{(D) }(xy)^{-1}$. ~\href{https://artofproblemsolving.com/wiki/index.php?title=User:Jiang147369}{jiang147369}

 


 


	\item (\textit{AHSME 1977}) If $t = \frac{1}{1 - \sqrt[4]{2}}$, then $t$ equals 


 $\text{(A)}\ (1-\sqrt[4]{2})(2-\sqrt{2})\qquad \text{(B)}\ (1-\sqrt[4]{2})(1+\sqrt{2})\qquad \text{(C)}\ (1+\sqrt[4]{2})(1-\sqrt{2}) \qquad \\ \text{(D)}\ (1+\sqrt[4]{2})(1+\sqrt{2})\qquad \text{(E)}-(1+\sqrt[4]{2})(1+\sqrt{2})$


 



 
 \textbf{Answer: }E



\textbf{Solution}

Solution by e\_power\_pi\_times\_i

 $t = \dfrac{1}{1-\sqrt[4]{2}} = (\dfrac{1}{1-\sqrt[4]{2}})(\dfrac{1+\sqrt[4]{2}}{1+\sqrt[4]{2}}) = (\dfrac{1+\sqrt[4]{2}}{1-\sqrt{2}})(\dfrac{1+\sqrt{2}}{1+\sqrt{2}}) = \dfrac{(1+\sqrt[4]{2})(1+\sqrt{2})}{-1} = \text{(E)}-(1+\sqrt[4]{2})(1+\sqrt{2})$.

 


	\item (\textit{AHSME 1977}) For every triple $(a,b,c)$ of non-zero real numbers, form the number $\frac{a}{|a|}+\frac{b}{|b|}+\frac{c}{|c|}+\frac{abc}{|abc|}$. 
The set of all numbers formed is


 $\textbf{(A)}\ {0} \qquad \textbf{(B)}\ \{-4,0,4\} \qquad \textbf{(C)}\ \{-4,-2,0,2,4\} \qquad \textbf{(D)}\ \{-4,-2,2,4\}\qquad \textbf{(E)}\ \text{none of these}$


 
 \textbf{Answer: }B



\textbf{Solution}

Solution by e\_power\_pi\_times\_i

 $\dfrac{x}{|x|} = 1$ or $-1$ depending whether $x$ is positive or negative. If $a$, $b$, and $c$ are positive, then the entire thing amounts to $4$. If one of the three is negative and the other two positive, the answer is $0$. If two of the three is negative and one is positive, the answer is $0$. If all three are negative, the answer is $-4$. Therefore the set is $\boxed{\textbf{(B)}\ \{-4,0,4\}}$.

 This is very similar to \href{https://artofproblemsolving.com/wiki/index.php/2017\_AMC\_8\_Problems/Problem\_21}{https://artofproblemsolving.com/wiki/index.php/2017\_AMC\_8\_Problems/Problem\_21}

 


	\item (\textit{AHSME 1977}) 
                    \begin{figure}[H]
                        centering
                        \includegraphics[width=0.5\textwidth]{Problem/AHSME_I_1977_9_Problem_Figure_1.png}
                    \end{figure}
                    


 

In the adjoining figure $\measuredangle E=40^\circ$ and arc $AB$, arc $BC$, and arc $CD$ all have equal length. Find the measure of $\measuredangle ACD$.


 $\textbf{(A) }10^\circ\qquad \textbf{(B) }15^\circ\qquad \textbf{(C) }20^\circ\qquad \textbf{(D) }\left(\frac{45}{2}\right)^\circ\qquad \textbf{(E) }30^\circ$


 
 \textbf{Answer: }B



\textbf{Solution}

Solution by e\_power\_pi\_times\_i

 If arcs $AB$, $BC$, and $CD$ are congruent, then $\measuredangle ACB = \measuredangle BDC = \measuredangle CBD = \theta$. Because $ABCD$ is cyclic, $\measuredangle CAD = \measuredangle CBD = \theta$, and $\measuredangle ADB = \measuredangle ACB = \theta$. Then, $\measuredangle EAD = \measuredangle EDA = \dfrac{180^\circ - 40^\circ}{2} = 70^\circ$. $\theta = 55^\circ$. $\measuredangle ACD = 180^\circ - 55^\circ - 110^\circ = \boxed{\textbf{(B) }15^\circ}$.

 


	\item (\textit{AHSME 1977}) If $(3x-1)^7 = a_7x^7 + a_6x^6 + \cdots + a_0$, then $a_7 + a_6 + \cdots + a_0$ equals 


 $\text{(A)}\ 0 \qquad  \text{(B)}\ 1 \qquad  \text{(C)}\ 64 \qquad  \text{(D)}\ -64 \qquad  \text{(E)}\ 128$


 
 \textbf{Answer: }E



\textbf{Solution}

Solution by e\_power\_pi\_times\_i

 Notice that if $x=1$, then $a_7x^7 + a_6x^6 + \cdots + a_0 =  a_7 + a_6 + \cdots + a_0$. Therefore the answer is $(3(1)-1)^7) = \boxed{\text{(E)}\ 128}$.

 


	\item (\textit{AHSME 1977}) For each real number $x$, let $\textbf{[}x\textbf{]}$ be the largest integer not exceeding $x$ 
(i.e., the integer $n$ such that $n\le x<n+1$). Which of the following statements is (are) true?


 $\textbf{I. [}x+1\textbf{]}=\textbf{[}x\textbf{]}+1\text{ for all }x \\ \textbf{II. [}x+y\textbf{]}=\textbf{[}x\textbf{]}+\textbf{[}y\textbf{]}\text{ for all }x\text{ and }y \\ \textbf{III. [}xy\textbf{]}=\textbf{[}x\textbf{]}\textbf{[}y\textbf{]}\text{ for all }x\text{ and }y$


 

$\textbf{(A) }\text{none}\qquad \textbf{(B) }\textbf{I }\text{only}\qquad \textbf{(C) }\textbf{I}\text{ and }\textbf{II}\text{ only}\qquad \textbf{(D) }\textbf{III }\text{only}\qquad \textbf{(E) }\text{all}$


 
 \textbf{Answer: }B



\textbf{Solution}

Solution by e\_power\_pi\_times\_i

 Notice that $\textbf{[}x\textbf{]}$ is just $\left \lfloor x \right \rfloor$. We see that $\textbf{I}$ is true, as adding by one does not change the fraction part of the number. Similarly, $\textbf{II}$ is false, because $\textbf{[}x+y\textbf{]}$ does not always equal $\textbf{[}x\textbf{]}+\textbf{[}y\textbf{]}$ (If both numbers were $1.5$, $\textbf{[}1.5+1.5\textbf{]} = 3$, $\textbf{[}1.5\textbf{]}+\textbf{[}1.5\textbf{]} = 2$). Because $\textbf{I}$ is true and $\textbf{II}$ is false, the answer is $\boxed{\textbf{(B) }\textbf{I }\text{only}}$.

 


	\item (\textit{AHSME 1977}) Al's age is $16$ more than the sum of Bob's age and Carl's age, and the square of Al's age is $1632$ more than the square of the sum of 
Bob's age and Carl's age. What is the sum of the ages of Al, Bob, and Carl?     


 $\text{(A)}\ 64 \qquad  \text{(B)}\ 94 \qquad  \text{(C)}\ 96 \qquad  \text{(D)}\ 102 \qquad  \text{(E)}\ 140$


 
 \textbf{Answer: }D



\textbf{Solution}

Solution by e\_power\_pi\_times\_i

 Denote Al's age, Bob's age, and Carl's age by $a$, $b$, and $c$, respectively. Then, $a = 16 + b + c$ and $a^2 = 1632 + b^2 + c^2$. Substituting the first equation into the second, $(16 + b + c)^2 = b^2 + c^2 + 2bc + 32b + 32c + 256 = b^2 + c^2 + 1632$. Thus, $bc + 16b + 16c = 688$, and $(b+16)(c+16) = 944$. Since $944 = 2^4\cdot59$, $(b,c) = (0,43)$ or $(43,0)$. Then $a + b + c = 2b + 2c + 16 = \boxed{\textbf{(D)}\ 102}$.

 


	\item (\textit{AHSME 1977}) If $a_1,a_2,a_3,\dots$ is a sequence of positive numbers such that $a_{n+2}=a_na_{n+1}$ for all positive integers $n$, 
then the sequence $a_1,a_2,a_3,\dots$ is a geometric progression


 $\textbf{(A) }\text{for all positive values of }a_1\text{ and }a_2\qquad\\ \textbf{(B) }\text{if and only if }a_1=a_2\qquad\\ \textbf{(C) }\text{if and only if }a_1=1\qquad\\ \textbf{(D) }\text{if and only if }a_2=1\qquad\\ \textbf{(E) }\text{if and only if }a_1=a_2=1$


 
 \textbf{Answer: }E



\textbf{Solution}

Solution by e\_power\_pi\_times\_i

 The first few terms are $a_1,a_2,a_1a_2,a_1a_2^2,a_1^2a_2^3,\dots$ . If this is a geometric progression, $\dfrac{a_2}{a_1} = a_1 = a_2 = a_1a_2$. $a_1=0,1$, $a_2=0,1$. Since this is a sequence of positive integers, then the answer must be $\boxed{\textbf{(E) }\text{if and only if }a_1=a_2=1 }$.

 


	\item (\textit{AHSME 1977}) How many pairs $(m,n)$ of integers satisfy the equation $m+n=mn$?


 $\textbf{(A) }1\qquad \textbf{(B) }2\qquad \textbf{(C) }3\qquad \textbf{(D) }4\qquad  \textbf{(E) }\text{more than }4$


 
 \textbf{Answer: }B



\textbf{Solution}

Solution by e\_power\_pi\_times\_i

 If $m+n=mn$, $mn-m-n = (m-1)(n-1)-1 = 0$. Then $(m-1)(n-1) = 1$, and $(m,n) = (2,2) or (0,0)$. The answer is $\boxed{\textbf{(B) }2}$.

 


	\item (\textit{AHSME 1977}) 
                    \begin{figure}[H]
                        centering
                        \includegraphics[width=0.5\textwidth]{Problem/AHSME_I_1977_15_Problem_Figure_1.png}
                    \end{figure}
                    


 Each of the three circles in the adjoining figure is externally tangent to the other two, and each side of the triangle is tangent to two of the circles. If each circle has radius three, then the perimeter of the triangle is


 $\textbf{(A) }36+9\sqrt{2}\qquad \textbf{(B) }36+6\sqrt{3}\qquad \textbf{(C) }36+9\sqrt{3}\qquad \textbf{(D) }18+18\sqrt{3}\qquad \textbf{(E) }45$


 
 \textbf{Answer: }D



\textbf{Solution}

Solution by e\_power\_pi\_times\_i

 Draw perpendicular lines from their if goe[ radii of the circles to the sides of the triangle, and lines from the radii of the circles to the vertices of the triangle. Because the triangle is equilateral, the lines divide the big triangle into a small triangle, three rectangles, and six small $30-60-90$ triangles. The Longer length of the rectangles is 6 because it is the radius of one circle + the radius of another ($3+3$). The length of one side is $2(3\sqrt{3})+6 = 6\sqrt{3}+6$. The perimeter is $\boxed{\textbf{(D) }18\sqrt{3}+18}$.

 


	\item (\textit{AHSME 1977}) If $i^2 = -1$, then the sum 
 \[\cos{45^\circ} + i\cos{135^\circ} + \cdots + i^n\cos{(45 + 90n)^\circ}  + \cdots + i^{40}\cos{3645^\circ}\]
equals 


 $\text{(A)}\ \frac{\sqrt{2}}{2} \qquad  \text{(B)}\ -10i\sqrt{2} \qquad  \text{(C)}\ \frac{21\sqrt{2}}{2} \qquad\\ \text{(D)}\ \frac{\sqrt{2}}{2}(21 - 20i) \qquad  \text{(E)}\ \frac{\sqrt{2}}{2}(21 + 20i)$


 
 \textbf{Answer: }D



\textbf{Solution}

Solution by e\_power\_pi\_times\_i

 Notice that the sequence repeats every $4$ terms. These $4$ terms are $\cos(45^\circ)$, $i\cos(135^\circ)$, $-\cos(225^\circ)$, and $-i\cos(315^\circ)$. This is because $\cos(\theta) = \cos(360^\circ - \theta)$. Also $\cos(\theta) = -\cos(180^\circ - \theta)$, so the $4$ terms are just $\cos(45^\circ)$, $-i\cos(45^\circ)$, $\cos(45^\circ)$, and $-i\cos(45^\circ)$, which adds up to $\sqrt{2}-i\sqrt{2}$. The sequence repeats $10$ times, and an extra $i^{40}\cos(3645^\circ)$. The sum is $10(\sqrt{2}-i\sqrt{2})+\dfrac{\sqrt{2}}{2} = \dfrac{21\sqrt{2}}{2} - 10i\sqrt{2} = \boxed{\text{(D) }\dfrac{\sqrt{2}}{2}(21-20i)}$.

 


	\item (\textit{AHSME 1977}) Three fair dice are tossed at random (i.e., all faces have the same probability of coming up). 
What is the probability that the three numbers turned up can be arranged to form an arithmetic progression with common difference one?


 $\textbf{(A) }\frac{1}{6}\qquad \textbf{(B) }\frac{1}{9}\qquad \textbf{(C) }\frac{1}{27}\qquad \textbf{(D) }\frac{1}{54}\qquad \textbf{(E) }\frac{7}{36}$


 
 \textbf{Answer: }B



\textbf{Solution}

Solution by e\_power\_pi\_times\_i

 
The denominator of the probability is $6^3 = 216$. We are asked to find the probability that the three numbers are consecutive. The three numbers can be arranged in $6$ ways, and there are $4$ triplets of consecutive numbers from $1-6$. The probability is $\dfrac{6\cdot4}{216} = \dfrac{24}{216} = \boxed{\text{(B) }\dfrac{1}{9}}$.

 


	\item (\textit{AHSME 1977}) If $y=(\log_23)(\log_34)\cdots(\log_n[n+1])\cdots(\log_{31}32)$ then 


 $\textbf{(A) }4<y<5\qquad \textbf{(B) }y=5\qquad \textbf{(C) }5<y<6\qquad \textbf{(D) }y=6\qquad \\ \textbf{(E) }6<y<7$


 
 \textbf{Answer: }B



\textbf{Solution}

Solution by e\_power\_pi\_times\_i

 
Note that $\log_{a}b = \dfrac{\log{b}}{\log{a}}$. Then $y=(\dfrac{\log3}{\log2})(\dfrac{\log4}{\log3})\cdots(\dfrac{\log32}{\log31}) = \dfrac{\log32}{\log2} = \log_232 = \boxed{\text{(B) }y=5}$.

 


	\item (\textit{AHSME 1977}) Let $E$ be the point of intersection of the diagonals of convex quadrilateral $ABCD$, and let $P,Q,R$, and $S$ be the centers of the circles 
circumscribing triangles $ABE, BCE, CDE$, and $ADE$, respectively. Then


 $\textbf{(A) }PQRS\text{ is a parallelogram}\\ \textbf{(B) }PQRS\text{ is a parallelogram if an only if }ABCD\text{ is a rhombus}\\ \textbf{(C) }PQRS\text{ is a parallelogram if an only if }ABCD\text{ is a rectangle}\\ \textbf{(D) }PQRS\text{ is a parallelogram if an only if }ABCD\text{ is a parallelogram}\\ \textbf{(E) }\text{none of the above are true}$


 
 \textbf{Answer: }A



\textbf{Solution 1}

Let $l_1$, $l_2$, $l_3$, and $l_4$ be the perpendicular bisectors of $AE$, $BE$, $CE$, and $DE$ respectively.  Then $P = l_1 \cap l_2$, $Q = l_2 \cap l_3$, $R = l_3 \cap l_4$, and $S = l_4 \cap l_1$.  Thus $l_1 = \overleftrightarrow{PS}$, $l_2 = \overleftrightarrow{PQ}$, $l_3 = \overleftrightarrow{QR}$, and $l_4 = \overleftrightarrow{RS}$.  

 Since $\overleftrightarrow{PS}$ and $\overleftrightarrow{QR}$ are both perpendicular to $\overleftrightarrow{AC}$,  $\overleftrightarrow{PS} \parallel \overleftrightarrow{QR}$.  Also, since $\overleftrightarrow{PQ}$ and $\overleftrightarrow{RS}$ are both perpendicular to $\overleftrightarrow{BD}$,  $\overleftrightarrow{PQ} \parallel \overleftrightarrow{RS}$.

 Because $PQRS$ has two pairs of parallel opposite sides, $\boxed{\textbf{(A) }PQRS\text{ is a parallelogram}}$.

 -j314andrews

 



\textbf{Solution 2}

Let $L$, $M$, $N$, and $O$ be the intersections of $AC$ and $PS$, $BD$ and $PQ$, $CA$ and $QR$, and $DB$ and $RS$ respectively. Since the \href{https://artofproblemsolving.com/wiki/index.php?title=Circumcenter}{circumcenter} of a triangle is determined by the intersection of its perpendicular bisectors, each of the pairs mentioned earlier are perpendicular to each other. Let $\angle SPQ=\alpha$. Then $\angle LEM=180^\circ-\alpha$ since $\angle PLC=PMD=90^\circ$. This means $\angle DEC=180^\circ-\alpha$, which similarly implies $\angle QRS=\alpha$. Through a similar method, we find that $\angle PSR=\angle PQR=180^\circ-\alpha$.

 Since the opposite angles of $PQRS$ are congruent, $PQRS$ is indeed a parallelogram - and we showed this without any assumption about the type of quadrilateral $ABCD$ was. Therefore, $PQRS$ is always a parallelogram, and so our answer is $\fbox{A}$. 

 


	\item (\textit{AHSME 1977}) $\begin{tabular}{ccccccccccccc}& & & & & & C & & & & & &\\  & & & & & C & O & C & & & & &\\  & & & & C & O & N & O & C & & & &\\  & & & C & O & N & T & N & O & C & & &\\  & & C & O & N & T & E & T & N & O & C & &\\  & C & O & N & T & E & S & E & T & N & O & C &\\  C & O & N & T & E & S & T & S & E & T & N & O & C \end{tabular}$


 

For how many paths consisting of a sequence of horizontal and/or vertical line segments, with each segment connecting a pair of adjacent letters in the diagram above, is the word CONTEST spelled out as the path is traversed from beginning to end?


 $\textbf{(A) }63\qquad \textbf{(B) }128\qquad \textbf{(C) }129\qquad \textbf{(D) }255\qquad  \textbf{(E) }\text{none of these}$


 
 \textbf{Answer: }E



\textbf{Solution}

The path will either start on the left side, start on the right side, or go down the middle. Say it starts on the left side. Instead of counting it starting from the $C$, let's start at the T and work backwards. Starting at the $T$, there are $2$ options to choose the $S$. Then, there are $2$ options to chose the $E$, $T$, $N$, $O$, and $C$, for a total of $2^6=64$ possibilities. But we need to subtract the case where it does down the middle, which makes it $63$ possibilities. The right side also has $63$ possibilities by symmetry. There is $1$ way for it to go down the middle, for a total of $63+63+1=127$ paths, so the answer is $\textbf{(E) }\text{none of these}$.

 


	\item (\textit{AHSME 1977}) For how many values of the coefficient a do the equations 
 \begin{align*}x^2+ax+1=0 \\ x^2-x-a=0\end{align*}
 have a common real solution?


 $\textbf{(A)}\ 0 \qquad \textbf{(B)}\ 1 \qquad \textbf{(C)}\ 2 \qquad \textbf{(D)}\ 3 \qquad \textbf{(E)}\ \infty$


 
 \textbf{Answer: }B



\textbf{Solution 1}

Subtracting the equations, we get $ax+x+1+a=0$, or $(x+1)(a+1)=0$, so $x=-1$ or $a=-1$. If $x=-1$, then $a=2$, which satisfies the condition. If $a=-1$, then $x$ is nonreal. This means that $a=-1$ is the only number that works, so our answer is $(B)$.

 



\textbf{Solution 2}

In the second equation, isolating $a$ yields $a = x^2 - x$.  Substituting into the first equation yields $x^2+(x^2-x)x+1=0$, which simplifies to $x^3+1=0$.  The only real root of $x^3+1$ is $-1$, and substituting $x=-1$ reduces both the original equations to $2-a=0$.  So $(x,a)=(-1,2)$ is the only real solution to these equations, and thus there is $\boxed{\textbf{(B)}\ 1}$ value of $a$ such that these equations have a common solution.

 -j314andrews

 


	\item (\textit{AHSME 1977}) If $f(x)$ is a real valued function of the real variable $x$, and $f(x)$ is not identically zero, 
and for all $a$ and $b$ $f(a+b)+f(a-b)=2f(a)+2f(b)$, then for all $x$ and $y$


 $\textbf{(A) }f(0)=1\qquad \textbf{(B) }f(-x)=-f(x)\qquad \textbf{(C) }f(-x)=f(x)\qquad \\ \textbf{(D) }f(x+y)=f(x)+f(y) \qquad \\ \textbf{(E) }\text{there is a positive real number }T\text{ such that }f(x+T)=f(x)$


 
 \textbf{Answer: }C



\textbf{Solution}

Substituting $(a,b)=(0,0)$ yields $2f(0)=4f(0)$, so $f(0)=0$.

 Substituting $(a,b)=(0,x)$ yields $f(x)+f(-x)=2f(-x)$ so $\boxed{\textbf{(C)}\ f(x)=f(-x)}$ for all $x$.

 Note: $f(x) = x^2$ satisfies the original equation and may be used to show that the other answer choices are incorrect.

 


	\item (\textit{AHSME 1977}) If the solutions of the equation $x^2+px+q=0$ are the cubes of the solutions of the equation $x^2+mx+n=0$, then


 $\textbf{(A) }p=m^3+3mn\qquad \textbf{(B) }p=m^3-3mn\qquad \textbf{(C) }p+q=m^3\qquad\\ \textbf{(D) }\left(\frac{m}{n}\right)^2=\frac{p}{q}\qquad  \textbf{(E) }\text{none of these}$


 
 \textbf{Answer: }B



\textbf{Solution}

Let $r_1$ and $r_2$ be the roots of the equation $x^2+mx+n=0$. Then, $r_1^3$ and $r_2^3$ are the roots of $x^2+px+q=0$. Applying \href{https://artofproblemsolving.com/wiki/index.php?title=Vieta\%27s\_Formulas}{Vieta's Formulas} to the first equation, we have $r_1 + r_2 = -m$ and $r_1r_2 = n$. Similarly, from the second equation, we have $r_1^3 + r_2^3 = -p$ and $r_1^3r_2^3 = q$. Cubing $r_1 + r_2 = -m$, we get $r_1^3 + r_2^3 + 3r_1r_2(r_1+r_2) = -m^3$. Plugging in from the other equations gives $-p + 3n(-m)=-m^3$. Rearranging we have $\boxed{\textbf{(B)} p=m^3-3mn.}$

 


	\item (\textit{AHSME 1977}) Find the sum $\frac{1}{1(3)}+\frac{1}{3(5)}+\dots+\frac{1}{(2n-1)(2n+1)}+\dots+\frac{1}{255(257)}$.


 $\textbf{(A) }\frac{127}{255}\qquad \textbf{(B) }\frac{128}{255}\qquad \textbf{(C) }\frac{1}{2}\qquad \textbf{(D) }\frac{128}{257}\qquad \textbf{(E) }\frac{129}{257}$


 
 \textbf{Answer: }D



\textbf{Solution}

Note that $\frac{1}{(2n-1)(2n+1)} = \frac{1/(2n-1)-1/(2n+1)}{2}$. Indeed, we find the series telescopes and is equal to $\frac{1-\frac{1}{257}}{2}$, which is evidently $\boxed {\frac{128}{257}}$

 


	\item (\textit{AHSME 1977}) Determine the largest positive integer n such that $1005!$ is divisible by $10^n$.


 $\textbf{(A) }102\qquad \textbf{(B) }112\qquad \textbf{(C) }249\qquad \textbf{(D) }502\qquad  \textbf{(E) }\text{none of these}$


 
 \textbf{Answer: }E



\textbf{Solution}

We first observe that since there will be more 2s than 5s in $1005!$, we are looking for the largest $n$ such that $5^n$ divides $1005!$. We will use the fact that:

 
 \[n = \left \lfloor {\frac{1005}{5^1}}\right \rfloor + \left \lfloor {\frac{1005}{5^2}}\right \rfloor + \left \lfloor {\frac{1005}{5^3}}\right \rfloor \cdots\]

 (This is an application of Legendre's formula).

 From $k=5$ and onwards, $\left \lfloor {\frac{1005}{5^k}}\right \rfloor = 0$. Thus, our calculation becomes

 
 \[n = \left \lfloor {\frac{1005}{5^1}}\right \rfloor + \left \lfloor {\frac{1005}{5^2}}\right \rfloor + \left \lfloor {\frac{1005}{5^3}}\right \rfloor + \left \lfloor {\frac{1005}{5^4}}\right \rfloor\]

 
 \[n = \left \lfloor {\frac{1005}{5}}\right \rfloor + \left \lfloor {\frac{1005}{25}}\right \rfloor + \left \lfloor {\frac{1005}{125}}\right \rfloor + \left \lfloor {\frac{1005}{625}}\right \rfloor\]

 
 \[n = 201 + 40 + 8 + 1 = \boxed{250}\]

 Since none of the answer choices equal 250, the answer is $\boxed{\textbf{(E)}}$.

 - mako17

 


	\item (\textit{AHSME 1977}) Let $a,b,c$, and $d$ be the lengths of sides $MN,NP,PQ$, and $QM$, respectively, of quadrilateral $MNPQ$. If $A$ is the area of $MNPQ$, then


 $\textbf{(A) }A=\left(\frac{a+c}{2}\right)\left(\frac{b+d}{2}\right)\text{ if and only if }MNPQ\text{ is convex}\\ \textbf{(B) }A=\left(\frac{a+c}{2}\right)\left(\frac{b+d}{2}\right)\text{ if and only if }MNPQ\text{ is a rectangle}\\ \textbf{(C) }A\le\left(\frac{a+c}{2}\right)\left(\frac{b+d}{2}\right)\text{ if and only if }MNPQ\text{ is a rectangle}\\ \textbf{(D) }A\le\left(\frac{a+c}{2}\right)\left(\frac{b+d}{2}\right)\text{ if and only if }MNPQ\text{ is a parallelogram}\\ \textbf{(E) }A\ge\left(\frac{a+c}{2}\right)\left(\frac{b+d}{2}\right)\text{ if and only if }MNPQ\text{ is a parallelogram}$


 
 \textbf{Answer: }B


	\item (\textit{AHSME 1977}) There are two spherical balls of different sizes lying in two corners of a rectangular room, each touching two walls and the floor. 
If there is a point on each ball which is $5$ inches from each wall which that ball touches and $10$ inches from the floor, 
then the sum of the diameters of the balls is


 $\textbf{(A) }20\text{ inches}\qquad \textbf{(B) }30\text{ inches}\qquad \textbf{(C) }40\text{ inches}\qquad \textbf{(D) }60\text{ inches}\qquad \\ \textbf{(E) }\text{not determined by the given information}$


 
 \textbf{Answer: }C



\textbf{Solution}

Let $r$ be the radius of one of the balls, in inches, $O$ be the center of the ball, and $A$ be the point which is $5$ inches away from each wall and $10$ inches above the floor.  Define a three-dimensional Cartesian coordinate system such that a corner of the room is the origin, the floor is in the positive $xy$-plane, one wall is in the positive $yz$-plane, and another wall is in the positive $xz$-plane.

 Then $O$ is at $(r, r, r)$ and $A$ is at $(5, 5, 10)$.  By the distance formula, $OA^2 = (r-5)^2 + (r-5)^2 + (r-10)^2 = r^2$.  This can be simplified to $2(r-15)(r-5) = 0$, which has solutions $r = 15$ and $r  = 5$.  So one ball has diameter $30$ inches while the other has diameter $10$ inches, thus the sum is $\boxed{\textbf{(C)}\ 40}$ inches.

 -j314andrews

 


	\item (\textit{AHSME 1977}) Let $g(x)=x^5+x^4+x^3+x^2+x+1$. What is the remainder when the polynomial $g(x^{12})$ is divided by the polynomial $g(x)$?


 $\textbf{(A) }6\qquad \textbf{(B) }5-x\qquad \textbf{(C) }4-x+x^2\qquad \textbf{(D) }3-x+x^2-x^3\qquad \\ \textbf{(E) }2-x+x^2-x^3+x^4$


 
 \textbf{Answer: }A



\textbf{Solution 1}

Let $r(x)$ be the remainder when $g(x^{12})$ is divided by $g(x)$. Then $r(x)$ is the unique polynomial such that
 \[g(x^{12}) - r(x) = x^{60} + x^{48} + x^{36} + x^{24} + x^{12} + 1 - r(x)\]is divisible by $g(x) = x^5 + x^4 + x^3 + x^2 + x + 1$, and $\deg r(x) < 5$.

 Note that $(x - 1)(x^5 + x^4 + x^3 + x^2 + 1) = x^6 - 1$ is a multiple of $g(x)$. Also,
 \[g(x^{12}) - 6 = x^{60} + x^{48} + x^{36} + x^{24} + x^{12} - 5 \\ = (x^{60} - 1) + (x^{48} - 1) + (x^{36} - 1) + (x^{24} - 1) + (x^{12} - 1).\]Each term is a multiple of $x^6 - 1$. For example,
 \[x^{60} - 1 = (x^6 - 1)(x^{54} + x^{48} + \dots + x^6 + 1).\]Hence, $g(x^{12}) - 6$ is a multiple of $x^6 - 1$, which means that $g(x^{12}) - 6$ is a multiple of $g(x)$. Therefore, the remainder is $\boxed{6}$. The answer is (A).

 



\textbf{Solution 2}

We express the quotient and remainder as follows.
 \[g(x^{12}) = Q(x) g(x) + R(x)\]Note that the solutions to $g(x)$ correspond to the 6th roots of unity, excluding $1$. Hence, we have $x^6 = 1$, allowing us to set:
 \[g(x^{12}) = 6\]
 \[g(x) = 0\]We have $5$ values of $x$ that return $R(x) = 6$. However, $g(x)$ is quintic, implying the remainder is of degree at most $4$. Since there are $5$ solutions, the only possibility is that the remainder is a constant $\boxed{6}$.

 



\textbf{Solution 3}

We can use the Chinese remainder theorem over $\mathbb{Q}[x].$ Since $g(x)=(x+1)(x^2-x+1)(x^2+x+1)$, $\mathbb{Q}[x]/g\cong \mathbb{Q}[x]/(x+1)\times \mathbb{Q}[x]/(x^2-x+1)\times \mathbb{Q}[x]/(x^2+x+1).$ This means that if we can find the remainder of $g(x^{12})$ modulo $x+1,x^2-x+1,x^2+x+1$, we can reconstruct the remainder modulo $g.$ We can further use that each factor is irreducible and that if $p(x)$ is an irreducible polynomial over $\mathbb{Q}$ with root $\alpha$, $\mathbb{Q}[x]/p\cong \mathbb{Q}(\alpha)$ so to evaluate the remainders of $g(x^{12})$, we just need to evaluate it on one of the roots of the irreducible factors. The first factor has root $-1$, the second has roots the primitive sixth roots of unity, and the third as roots the primitive cube roots of unity (this is easily seen as $g(x)(x-1)=x^6-1$). Evaluating $g(x^{12})$ on each of these values yields $g(1)=6$ so the remainder is $6$ on each factor on the right of the isomorphism. Hence, by the Chinese remainder theorem, the remainder modulo $g$ must be $\boxed{6}$ as well.

 



\textbf{Solution 4 (Remainder Theorem)}

First, we factor \(g(x)\) as \((x^3+1)(x^2+x+1)\). Substituting to find \(g(x^{12})\), we have  

 
 \[g(x^{12})=(x^{36}+1)(x^{24}+x^{12}+1).\]

 We want the remainder of  

 
 \[\frac{g(x^{12})}{g(x)}=\frac{(x^{36}+1)(x^{24}+x^{12}+1)}{(x^{3}+1)(x^{2}+x+1)}.\]

 Splitting it into two expressions, we can find the remainder of  
 \[\frac{x^{36}+1}{x^{3}+1}\quad \text{and} \quad\frac{x^{24}+x^{12}+1}{x^{2}+x+1},\]then multiply the results.  

 
                \begin{itemize}
                    \item First expression:
                \end{itemize}


                The roots of \(x^{3}+1\) are \(-1\), \(\text{cis}(60^\circ)\), and \(\text{cis}(-60^\circ)\). Plugging each of these roots into \(x^{36}+1\) yields \(2\). Since the remainder must have degree at most 2 and takes the constant value \(2\) at three distinct points, the remainder is \(r(x)=2\).  

 
                \begin{itemize}
                    \item Second expression:
                \end{itemize}


                The roots of \(x^{2}+x+1\) are \(\text{cis}(120^\circ)\) and \(\text{cis}(-120^\circ)\). Plugging each of these roots into \(x^{24}+x^{12}+1\) yields \(3\). Since the remainder must have degree at most 1 and takes the constant value \(3\) at two distinct points, the remainder is \(r(x)=3\).  

 Multiplying the two remainders, we get \(2\cdot 3=6\), or \(\boxed{\textbf{(A)}}\).~Voidling

 


	\item (\textit{AHSME 1977}) Find the smallest integer $n$ such that $(x^2+y^2+z^2)^2\le n(x^4+y^4+z^4)$ for all real numbers $x,y$, and $z$.


 $\textbf{(A) }2\qquad \textbf{(B) }3\qquad \textbf{(C) }4\qquad \textbf{(D) }6\qquad  \textbf{(E) }\text{There is no such integer n}$


 
 \textbf{Answer: }B



\textbf{Solution (Official MAA)}

Let $a = x^2$, $b = y^2$, $c = z^2$. Then 

 
 \[0 \leq (a-b)^2 + (b-c)^2 + (c-a)^2\]
 \[\dfrac{ab+bc+ca}{a^2 + b^2 + c^2} \leq 1;\]
 \[\dfrac{a^2 + b^2 + c^2 + 2(ab+bc+ca)}{a^2 + b^2 + c^2} \leq 3;\]
 \[(a+b+c)^2 \leq 3(a^2 + b^2 + c^2).\]

 Therefore $n\leq 3$. Choosing $a = b = c > 0$ shows $n$ is not less than three.

 


	\item (\textit{AHSME 1977}) 
                    \begin{figure}[H]
                        centering
                        \includegraphics[width=0.5\textwidth]{Problem/AHSME_I_1977_30_Problem_Figure_1.png}
                    \end{figure}
                    


 If $a,b$, and $d$ are the lengths of a side, a shortest diagonal and a longest diagonal, respectively, 
of a regular nonagon (see adjoining figure), then


 $\textbf{(A) }d=a+b\qquad \textbf{(B) }d^2=a^2+b^2\qquad \textbf{(C) }d^2=a^2+ab+b^2\qquad\\ \textbf{(D) }b=\frac{a+d}{2}\qquad  \textbf{(E) }b^2=ad$


 
 \textbf{Answer: }A



\textbf{Solution 1}

By the law of cosines we can get the following expressions for $d^{2}$ and $b^{2}$:

 
 \[d^{2}=2b^{2}(1-\cos (100^\circ))\] 
 \[b^{2}=2a^{2}(1-\cos (140^\circ))\]

 We can substitute what we got for $b^2$ into the expression for $d^{2}$:

 
 \[d^{2}=4a^{2}(1-\cos (140^\circ))(1-\cos (100^\circ))=4a^{2}(1-(\cos (100^\circ)+\cos (140^\circ))+\cos (100^\circ)\cos (140^\circ))\]

 Now apply sum-to-product and product-to-sum identities:
 \[d^{2}=4a^{2}(1-(2\cos (120^\circ)\cos (20^\circ))+\frac{1}{2}(\cos (240^\circ)+\cos (40^\circ)))\]Simplifying further gives us: 
 \[d^{2}=4a^{2}(1+\cos (20^\circ)-\frac{1}{4}+\frac{1}{2}\cos (40^\circ))\]After using the fact that $\cos (40^\circ)=2\cos^2 (20^\circ)-1$, it's not hard to see that the expression in the parentheses is equal to $(\cos (20^\circ)+\frac{1}{2})^{2}$. So we can square-root both sides to find the expression for $d$:

 
 \[d=2a(\cos (20^\circ)+\frac{1}{2})\]Now let's look at the expression for $b^2$. We can apply the reverse of the double angle identity to show that $1-\cos (140^\circ)$ equals $2\sin^2 (70^\circ))$. So if we square root the entire expression we get that 
 \[b=2a\sin (70^\circ)=2a\cos (20^\circ)\]We now have everything in terms of $a$. Luckily when we consider choice $\textbf{A}$ we can verify without much work that this must be the answer.

 Solution by harita19

 
                \begin{itemize}
                    \item NOTE: a much easier solution exists by drawing some lines and recognizing that the nonagon is cyclic, but for those of use who use algebra in every geometry problem, this is the best solution.
                \end{itemize}


                



\textbf{Solution 2}

Using Van Schooten's Theorem, we draw an equilateral triangle in the nonagon to see that $a+b=d$.

 


 



\textbf{Solution 3}

Reflect $b$ such that it connects the left vertex of the $a$ to the $2$nd vertex counterclockwise. This length is also $b$. Notice that when connecting the two endpoints of the $b$ segments, we achieve a length of $d$. We now have an isosceles trapezoid with bases $a$ and $d$. 

 Next, we can find the angle between the $a$ and the $b$ segments as follows. The interior angle of a regular nonagon measures $180-\frac{360}{9}=140^{\circ}$. Thus, the isosceles triangle formed from two of the side lengths of the nonagon and $b$ has angles $140$, $\frac{180-140}{2}$, and $\frac{180-140}{2}=20$. Then, the angle between $a$ and $b$ is $140-20=120$. We can therefore deduce that the other angles are $120$, $60$, and $60$ degrees.

 We can drop a perpendicular from both endpoints of segment $a$ to segment $d$ to get a rectangle and two $30$-$60$-$90$ triangles. Each of the triangles has a shorter leg length of $\frac{d-a}{2}$, and hence, we can deduce that $\frac{d-a}{2} \cdot 2=b$. Hence, the correct answer is $\mathbf{(A)} \; d=a+b$

 - Solution by scjh999999

 



\textbf{Solution 4}

Draw chords of length $a$, $b$, and $d$ to create two triangles as follows.

 
                    \begin{figure}[H]
                        centering
                        \includegraphics[width=0.5\textwidth]{Solution/AHSME_I_1977_30_Solution_Figure_1.png}
                    \end{figure}
                    

 Because each edge of a nonagon has a central angle of $40^{\circ}$, and inscribed angles are half the corresponding central angle, the two triangles thus created are equilateral.  This means that $d = a + b$.

 - Solution by Dr.B.

 


\end{enumerate}